% Theorem-facing replacement for Section 2.

\section{The complete independence-number phase}
\label{the-complete-independence-number-phase}

Put
\[
  q:=\log 2,
  \qquad L:=\log n,
  \qquad \ell:=\log\log n,
  \qquad C:=1+\log q-q.
\]
The usual independence-number center can then be written without a change of
logarithm base as
\begin{equation*}
  \alpha_0
  =2\log_2 n-2\log_2\log_2 n+2\log_2(\mathrm e/2)+1
  =\frac{2(L-\ell+C)}q+1.
  \tag{2.1}
\end{equation*}
Let
\[
  \alpha:=\lfloor\alpha_0\rfloor,
  \qquad
  \delta:=\alpha_0-\alpha,
  \qquad
  b:=1-\delta.
\]
Thus $0\le\delta<1$, $0<b\le1$, and
\[
  \alpha=\frac{2(L-\ell+C)}q+b.
\]
All estimates below are uniform for the closed parameter range
$0\le\delta\le1$; this includes integer sequences approaching either endpoint
of the actual half-open phase interval.

\begin{lemma}[Uniform phase expansion and adjacent-size control]
There is a bounded continuous function $K\colon[0,1]\to\mathbb R$, an
absolute constant $C_{\mathrm{ph}}>0$, and a deterministic sequence
\[
  \varepsilon_n^{\mathrm{ph}}
  :=C_{\mathrm{ph}}\frac{1+\ell^2}{L}
  \longrightarrow0
\]
such that, for every $0\le\delta\le1$,
\begin{equation*}
  \log\mu_\alpha
  =
  \delta L
  +\left(\frac2q-\frac12-\delta\right)\ell
  +K(\delta)
  +E_n(\delta),
  \qquad
  |E_n(\delta)|\le\varepsilon_n^{\mathrm{ph}}.
  \tag{2.2}
\end{equation*}
Moreover, uniformly in the same phase,
\begin{equation*}
\begin{split}
  \log\mu_{\alpha+2}
  &=(\delta-2)L
    +\left(\frac2q+\frac32-\delta\right)\ell+O(1),\\
  \mu_{\alpha+2}
  &\le \exp\{-L+C_2\ell\}=o(1)
\end{split}
  \tag{2.3}
\end{equation*}
for an absolute $C_2$, while, with $a:=\alpha-2$,
\begin{equation*}
  \mu_a
  \ge c\,n^2(\log n)^{2/q-5/2}
  \tag{2.4}
\end{equation*}
for an absolute $c>0$ and all sufficiently large $n$.
\end{lemma}

\begin{proof}
Because $\alpha=O(L)$ uniformly in the phase, one has $\alpha<n/2$ for all
sufficiently large $n$. For $0\le x\le1/2$,
$-2x\le\log(1-x)\le-x$. Hence
\[
  \log(n)_\alpha
  =\alpha L+\sum_{j=0}^{\alpha-1}\log\!\left(1-\frac jn\right)
  =\alpha L+R_{n,\alpha},
  \qquad
  |R_{n,\alpha}|\le\frac{\alpha^2}{n}.
\]
Using the uniform Stirling remainder
\[
  \log(\alpha!)
  =\alpha\log\alpha-\alpha
   +\frac12\log(2\pi\alpha)+O(1/\alpha)
\]
in
\[
  \log\mu_\alpha
  =\log(n)_\alpha-\log(\alpha!)-q\binom\alpha2
\]
gives
\begin{equation*}
  \log\mu_\alpha
  =
  \alpha\left(
    L-\log\alpha+1-\frac q2(\alpha-1)
  \right)
  -\frac12\log(2\pi\alpha)
  +O(1/L).
  \tag{2.5}
\end{equation*}
Every remainder here is independent of $\delta$.

Write
\[
  \alpha=\frac{2L}{q}(1+x_n),
  \qquad
  x_n:=\frac{-\ell+C+qb/2}{L}.
\]
Uniformly for $0\le b\le1$, one has $|x_n|\le1/2$ for all sufficiently
large $n$. The Taylor formula
$\log(1+x)=x+O(x^2)$ therefore gives
\[
  \log\alpha
  =\ell+q-\log q
   +\frac{-\ell+C+qb/2}{L}
   +O\!\left(\frac{\ell^2}{L^2}\right).
\]
Since
\[
  \frac q2(\alpha-1)=L-\ell+C-\frac{q\delta}{2},
\]
we obtain the uniform expansion
\begin{equation*}
  L-\log\alpha+1-\frac q2(\alpha-1)
  =
  \frac{q\delta}{2}
  +\frac{\ell-C-qb/2}{L}
  +O\!\left(\frac{\ell^2}{L^2}\right).
  \tag{2.6}
\end{equation*}
Multiplying by
$\alpha=2(L-\ell+C)/q+b$ and collecting terms yields
\[
\begin{split}
  \alpha\left(
    L-\log\alpha+1-\frac q2(\alpha-1)
  \right)
  ={}&
  \delta L+\left(\frac2q-\delta\right)\ell\\
  &+\delta C+\frac q2\delta(1-\delta)
    -\frac{2C}{q}-(1-\delta)
    +O\!\left(\frac{1+\ell^2}{L}\right).
\end{split}
\]
Also
\[
  -\frac12\log(2\pi\alpha)
  =-rac12\ell-rac12\log(2\pi)-\frac q2
    +\frac12\log q+O(\ell/L).
\]
Thus (2.2) holds with
\begin{equation*}
\begin{split}
  K(\delta):={}&
  \delta C+\frac q2\delta(1-\delta)-\frac{2C}{q}-(1-\delta)\\
  &-\frac12\log(2\pi)-\frac q2+\frac12\log q.
\end{split}
  \tag{2.7}
\end{equation*}
This function is continuous, hence bounded, on $[0,1]$. Enlarging one absolute
constant absorbs all displayed remainders into
$\varepsilon_n^{\mathrm{ph}}$.

The adjacent-size ratios are exact:
\begin{equation*}
  \frac{\mu_{s+1}}{\mu_s}
  =\frac{n-s}{s+1}2^{-s},
  \qquad
  \frac{\mu_{s-1}}{\mu_s}
  =\frac{s}{n-s+1}2^{s-1}.
  \tag{2.8}
\end{equation*}
Equation (2.1) gives the exact phase representation
\[
  2^\alpha
  =\exp(q\alpha)
  =\exp(2C+qb)\frac{n^2}{L^2}.
\]
Since $0\le b\le1$, the multiplicative factor $\exp(2C+qb)$ is bounded
above and below by positive absolute constants. Therefore, uniformly for
$s=\alpha+O(1)$,
\[
  \frac{\mu_{s+1}}{\mu_s}=\Theta(L/n),
  \qquad
  \frac{\mu_{s-1}}{\mu_s}=\Theta(n/L).
\]
In particular,
\begin{equation*}
  \mu_{\alpha+2}
  =\mu_\alpha\,\Theta(L^2/n^2),
  \qquad
  \mu_{\alpha-2}
  =\mu_\alpha\,\Theta(n^2/L^2).
  \tag{2.8a}
\end{equation*}
Taking logarithms and using (2.2) proves the first line of (2.3). Since
$\delta\le1$, its leading term is at most $-L$, and the coefficient of
$\ell$ is uniformly bounded; this proves the second line.

For the lower adjacent size, (2.2) and (2.8a) give
\[
\begin{split}
  \log\mu_{\alpha-2}
  ={}&2L+\left(\frac2q-\frac52\right)\ell
      +\delta(L-\ell)+K(\delta)+O(1).
\end{split}
\]
For large $n$, $L-\ell>0$, while $K$ is bounded below. Exponentiating gives
(2.4), with one absolute constant $c$ and one eventuality threshold valid for
the complete phase.
\end{proof}

Let $X_{\alpha+2}$ denote the number of independent sets of size
$\alpha+2$. Then
$\mathbb E X_{\alpha+2}=\mu_{\alpha+2}$, and the event
$\alpha(G_n)>\alpha+1$ implies $X_{\alpha+2}\ge1$. Consequently Markov's
inequality and (2.3) give the deterministic cap error
\begin{equation*}
  \mathbb P\bigl(\alpha(G_n)>\alpha+1\bigr)
  \le\mu_{\alpha+2}
  =:\varepsilon_n^{\mathrm{cap}}
  \longrightarrow0
  \tag{2.9}
\end{equation*}
uniformly along the full sequence of integers.